% !TeX root = ../../Book4.tex
% 纲要标题：### 7.1 链同伦
% 纲要标签：b4:sec:0601
% 纲要位置：计划纲要.md，第 3344 行。

\section{链同伦}
\label{b4:sec:0601}

同调看不见的复形映射未必都能用同伦消去。先研究可显式写成 \(\mathrm{d}h+h\mathrm{d}\) 的差别，便能区分零伦、可缩与仅仅正合。

% 纲要内容（仅作写作计划，尚非教材正文）：
%
% * 同伦定义与同调不变性
% * 零伦映射及可缩复形
% * 正合复形未必可缩
%

\subsection{同伦与同调不变性}
\label{b4:subsec:0601:01}

两个复形映射可能逐项不同，却在同调上完全相同。最容易控制的一类差别是能够写成“先作微分再作某个降次映射，加上先降次再作微分”。这类差别在余圈上自动变成余边。

\ProseParagraphBreak
\cref{b4:def:cochain-complex,b4:def:complex-map}在交换范畴中定义了复形与复形映射；在加性范畴中，同样可以取对象族 \(C^n\) 和满足 \(\mathrm{d}_C^{n+1}\mathrm{d}_C^n=0\) 的态射 \(\mathrm{d}_C^n:C^n\to C^{n+1}\)，并要求复形映射与微分交换。这些定义只用到零态射和复合。同伦还用到态射的加减法，因而也可在加性范畴中讨论；核、像及同调对象则仍在交换范畴中使用。

\begin{definition}\label{b4:def:chain-homotopy}
设 \(\mathcal A\) 是加性范畴，\(C,D\) 为 \(\mathcal A\) 中的上链复形，\(f,g:C\to D\) 是复形映射。若有态射族 \(h^n:C^n\to D^{n-1}\)（\(n\in\mathbb Z\)），使
\[
f^n-g^n=\mathrm{d}_D^{n-1}h^n+h^{n+1}\mathrm{d}_C^n
\quad(n\in\mathbb Z),
\]
则称 \(f\) 与 \(g\) \term[lian-tonglun]{链同伦}[chain homotopic]，记为 \(f\simeq g\)，并称 \(h\) 为一个链同伦。链指标中同一个公式写为 \(f_n-g_n=\partial_{D,n+1}h_n+h_{n-1}\partial_{C,n}\)，其中 \(h_n:C_n\to D_{n+1}\)。
\end{definition}
\symidx{\(f\simeq g\)：复形映射之间的链同伦关系}

上链指标下同伦降一次，链指标下升一次，两者都与微分方向相反。在次数 \(n\)，两个复合
\[
\mathrm{d}_D^{n-1}h^n:C^n\longrightarrow D^n,
\qquad h^{n+1}\mathrm{d}_C^n:C^n\longrightarrow D^n
\]
分别先降次再升次、先升次再降次，因而都能与 \(f^n-g^n\) 相加比较。这里保留“链同伦”这一通用名称，不因为指标升次而另换概念。

\ProseParagraphBreak
对模复形，\cref{b4:def:hom-complex}中的 Hom 复形恰好解释了这个公式的来源。态射族 \(h\) 是次数 \(-1\) 的元素，其微分为
\[
\mathrm{d}_{\operatorname{Hom}}h
=\mathrm{d}_Dh-(-1)^{-1}h\mathrm{d}_C
=\mathrm{d}_Dh+h\mathrm{d}_C.
\]
所以 \(f\simeq g\) 就是说零次余圈 \(f-g\) 是 Hom 复形中的余边。

\begin{proposition}\label{b4:prop:homotopy-cohomology}
设 \(\mathcal A\) 是加性范畴。在任意两个固定的 \(\mathcal A\) 中的上链复形之间，同伦是复形映射的等价关系，且与加法、左右复合相容。若 \(\mathcal A\) 还是交换范畴，\(f,g:C\to D\) 为同伦的复形映射，则对每个 \(n\in\mathbb Z\)，有 \(H^n(f)=H^n(g)\)。
\end{proposition}
\begin{proof}
零同伦、\(-h\) 及同伦之和分别给出自反、对称和传递性。若 \(u:C'\to C\)、\(v:D\to D'\) 为复形映射，则
\[
v(f-g)u=\mathrm{d}_{D'}(vhu)+(vhu)\mathrm{d}_{C'};
\]
这证明复合相容性。两个差的同伦相加给出加法相容性。

再设 \(\mathcal A\) 是交换范畴，记 \(z_C^n,z_D^n\) 为余圈的核嵌入，\(j_D^n:B^n(D)\to Z^n(D)\) 为余边到余圈的嵌入，并写
\[
\mathrm{d}_D^{n-1}=z_D^nj_D^n\pi_D^{n-1},
\qquad \pi_D^{n-1}:D^{n-1}\twoheadrightarrow B^n(D).
\]
设 \(f_Z^n,g_Z^n:Z^n(C)\to Z^n(D)\) 是 \(f,g\) 在余圈上诱导的态射。把同伦公式右复合 \(z_C^n\)，含 \(h^{n+1}\mathrm{d}_C^n\) 的项消失，再消去单态射 \(z_D^n\)，得到
\[
f_Z^n-g_Z^n=j_D^n\pi_D^{n-1}h^nz_C^n.
\]
同调商映射 \(q_D^n:Z^n(D)\to H^n(D)\) 满足 \(q_D^nj_D^n=0\)，故 \(q_D^n(f_Z^n-g_Z^n)=0\)。由 \(q_C^n:Z^n(C)\to H^n(C)\) 为满态射，得到 \(H^n(f)-H^n(g)=0\)。在模中，这就是余圈 \(z\) 的两个像之差 \(f(z)-g(z)=\mathrm{d}_Dh(z)\) 是余边。
\end{proof}

因此对任意加性范畴 \(\mathcal A\)，都可以把同伦的映射视为同一个态射，得到\term[tonglun-fanchou]{同伦范畴}[homotopy category] \(K(\mathcal A)\)。对象仍是复形，态射则是复形映射的同伦类；上述相容性保证态射的复合与加法不依赖代表。\symidx{\(K(\mathcal A)\)：复形的同伦范畴}例如投射模构成的加性范畴也有同伦范畴，不必要求它本身是交换范畴。

\ProseParagraphBreak
若 \(C,D\) 是左 \(R\)-模复形，则
\[
\begin{aligned}
Z^0\operatorname{Hom}_R^\bullet(C,D)
&=\{C\to D\text{ 的复形映射}\},\\
B^0\operatorname{Hom}_R^\bullet(C,D)
&=\{\mathrm{d}_Dh+h\mathrm{d}_C\mid h\text{ 的次数为 }-1\}.
\end{aligned}
\]
右边第二个集合恰是零伦映射所成的子群。将零次余圈模掉这个子群，便是将两个同伦的复形映射视为相同，故
\[
\operatorname{Hom}_{K(R\text{-}\mathbf{Mod})}(C,D)
=H^0\operatorname{Hom}_R^\bullet(C,D),
\]
这把同伦类的定义与 Hom 复形的零次同调联系起来。

\subsection{零伦映射及可缩复形}
\label{b4:subsec:0601:02}

\begin{definition}\label{b4:def:contractible-complex}
设 \(\mathcal A\) 为加性范畴，\(C,D\) 为其中的上链复形，\(f:C\to D\) 为复形映射。若 \(f\) 与零映射同伦，则称 \(f\) 为\term[linglun-yingshe]{零伦映射}[null-homotopic map]。若 \(1_C\simeq0\)，则称 \(C\) 为\term[kesuo-fuxing]{可缩复形}[contractible complex]；满足 \(1_{C^n}=\mathrm{d}_C^{n-1}h^n+h^{n+1}\mathrm{d}_C^n\) 的态射族 \(h^n:C^n\to C^{n-1}\)（\(n\in\mathbb Z\)）称为收缩同伦。若存在复形映射 \(g:D\to C\)，使 \(gf\simeq1_C\)、\(fg\simeq1_D\)，则称 \(f\) 为\term[tonglun-dengjia]{同伦等价}[homotopy equivalence]。
\end{definition}

同伦等价所需的 \(g\) 是同伦逆，两个复合只要求与恒等映射同伦；复形同构的严格逆当然满足这一要求，反向却不必成立。可缩复形在交换范畴中必零调，因为恒等同态在同调上等于零。反向所缺的条件是各次数的短正合列能否分裂：分裂把 \(C^n\) 分成当前的余圈与下一次余边的选定提升，微分便只负责把后一个部分送到下一项。

\begin{proposition}\label{b4:prop:contractible-split-exact}
设 \(\mathcal A\) 是交换范畴，\(C\) 为其中的上链复形。它可缩，当且仅当它零调，且对每个 \(n\in\mathbb Z\)，短正合列
\[
0\longrightarrow Z^n(C)\xrightarrow{z_C^n}C^n
\xrightarrow{\widetilde{\mathrm{d}}_C^n} Z^{n+1}(C)\longrightarrow0
\]
均分裂；其中 \(z_C^n:Z^n(C)\to C^n\) 为核态射，\(\widetilde{\mathrm{d}}_C^n\) 由 \(\mathrm{d}_C^n=z_C^{n+1}\widetilde{\mathrm{d}}_C^n\) 定义，并在零调假设下为满态射。
\end{proposition}
\begin{proof}
若 \(1_C=\mathrm{d}h+h\mathrm{d}\)，则已知 \(C\) 零调。\(\widetilde{\mathrm{d}}_C^n\) 是微分的像分解中到 \(Z^{n+1}(C)=B^{n+1}(C)\) 的满态射；这里满的是这个分量，不是到整个 \(C^{n+1}\) 的微分。收缩同伦 \(h^{n+1}\) 将下一次余圈降回 \(C^n\)，因而截面的候选是
\[
s^{n+1}\coloneqq h^{n+1}z_C^{n+1}:Z^{n+1}(C)\longrightarrow C^n.
\]
把次数 \(n+1\) 的同伦公式右复合 \(z_C^{n+1}\)，由 \(\mathrm{d}_C^{n+1}z_C^{n+1}=0\) 得
\[
z_C^{n+1}=\mathrm{d}_C^nh^{n+1}z_C^{n+1}
=z_C^{n+1}\widetilde{\mathrm{d}}_C^nh^{n+1}z_C^{n+1}.
\]
消去单态射 \(z_C^{n+1}\)，得到
\[
\widetilde{\mathrm{d}}_C^n\bigl(h^{n+1}z_C^{n+1}\bigr)
=1_{Z^{n+1}(C)}.
\]
所以 \(s^{n+1}\) 确是 \(\widetilde{\mathrm{d}}_C^n\) 的截面。

反过来，选各短列的截面 \(s^{n+1}:Z^{n+1}(C)\to C^n\)。分裂短正合列给出的同构具体为
\[
\alpha^n=(z_C^n,s^{n+1}):Z^n(C)\oplus Z^{n+1}(C)\xrightarrow{\sim}C^n.
\]
第一分量由微分消去，第二分量经微分成为下一项的余圈嵌入，故在这些双积坐标中，微分为 \(\mathrm{d}(z,w)=(w,0)\)。定义 \(h^n(z,w)=(0,z)\in Z^{n-1}(C)\oplus Z^n(C)\)，则 \(\mathrm{d}h(z,w)=(z,0)\)、\(h\mathrm{d}(z,w)=(0,w)\)，相加等于恒等。这里的坐标公式表示双积的分量态射，因而计算在一般交换范畴中也成立。
\end{proof}

例如在次数 \(n,n+1\) 放同一对象 \(M\)，微分为 \(1_M\)，其余各项为零。只取 \(h^{n+1}=1_M\)，其余同伦分量为零，就得到收缩同伦。它的两个非零项可以很大，同调却为零；可缩所说的“消去”是通过显式同伦实现的，并不是逐项零对象。

\subsection{正合复形的非可缩例子}
\label{b4:subsec:0601:03}

\begin{example}\label{b4:exm:acyclic-not-contractible}
把交换群短正合列 \(0\to\mathbb Z\xrightarrow{2}\mathbb Z\to\mathbb Z/2\mathbb Z\to0\) 放在次数 \(0,1,2\)，便得到零调复形。若它可缩，记次数一的商映射为 \(\mathrm{d}^1\)，由于 \(\mathrm{d}^2=0\)，次数二的同伦公式就是 \(1_{\mathbb Z/2\mathbb Z}=\mathrm{d}^1h^2\)。这要求 \(h^2:\mathbb Z/2\mathbb Z\to\mathbb Z\) 为商映射的截面。但任何这样的群同态都为零，因此不存在截面，复形不可缩。
\end{example}

\begin{example}\label{b4:exm:acyclic-free-periodic}
令 \(R=\mathbb Z/4\mathbb Z\)，在每个整数次数取自由模 \(R\)，微分均为乘 \(2\)。因为核与像同为 \(2R\)，复形正合。每个 \(R\)-线性 \(h^n:R\to R\) 都是乘某个 \(a_n\in R\)。若有收缩同伦，便须
\[
1=2a_n+2a_{n+1}\quad\text{在}\;R\;\text{中成立},
\]
右端是偶数剩余类，矛盾。这说明即使各项自由，无界正合复形也未必可缩。

\ProseParagraphBreak
此时余圈 \(2R\cong R/(2)\) 不是投射模。否则满射 \(R\xrightarrow{2}2R\) 会有截面 \(s:2R\to R\)；由 \(2s(2)=s(2\cdot2)=0\) 可知 \(s(2)\in2R\)，所以乘 \(2\) 又将它送到零，无法得到所要求的 \(2\)。因此各项自由并没有保证上述余圈短正合列分裂。
\end{example}

作为对照，对任意含幺环 \(R\)，各项均为投射幺性左 \(R\)-模的上有界零调复形必可缩。取 \(b\) 使 \(n>b\) 时 \(C^n=0\)，则 \(Z^b(C)=C^b\) 投射；于是 \(0\to Z^{b-1}(C)\to C^{b-1}\to Z^b(C)\to0\) 分裂，且 \(Z^{b-1}(C)\) 作为投射模的直和因子仍投射。如此向较低次数归纳，所有余圈短正合列都分裂，再用上述命题得到可缩性。周期反例没有可供这一步归纳开始的最高次数。后面比较投射消解时，也将从一个有限端开始逐次提升同伦，不能在没有起点的无界复形上直接照搬。
